% -*- TeX -*-
\documentclass[aspectratio=169]{beamer}

\title{PyLith v3.0 Tutorial}
\subtitle{Quasi-static Elasticity with Prescribed Fault Slip}
\author{Brad Aagaard\\
  Charles Williams \\
  Matthew Knepley}
\institute{\includegraphics[scale=1.5]{../../logos/cig_logo_dots}%
  \hspace{4em}%
\raisebox{1em}{\includegraphics[scale=1.0]{../../logos/cig_short_pylith}}}
\date{June 20, 2022}


% ---------------------------------------------------- CUSTOMIZATION
\usetheme{CIG}
% Style information for PyLith presentations.

% Colors
\definecolor{ltorange}{rgb}{1.0, 0.74, 0.41} % 255/188/105
\definecolor{orange}{rgb}{0.96, 0.50, 0.0} % 246/127/0

\definecolor{ltred}{rgb}{1.0, 0.25, 0.25} % 255/64/64
\definecolor{red}{rgb}{0.79, 0.00, 0.01} % 201/0/3

\definecolor{ltpurple}{rgb}{0.81, 0.57, 1.00} % 206/145/255
\definecolor{purple}{rgb}{0.38, 0.00, 0.68} % 97/1/175

\definecolor{ltblue}{rgb}{0.2, 0.73, 1.0} % 51/187/255
\definecolor{mdblue}{rgb}{0.28, 0.50, 0.80} % 72/128/205
\definecolor{blue}{rgb}{0.12, 0.43, 0.59} % 30/110/150

\definecolor{ltltgreen}{rgb}{0.7, 1.00, 0.7} % 96/204/14
\definecolor{ltgreen}{rgb}{0.37, 0.80, 0.05} % 96/204/14
\definecolor{green}{rgb}{0.23, 0.49, 0.03} % 59/125/8
  
\definecolor{dkslate}{rgb}{0.18, 0.21, 0.28} % 47/53/72
\definecolor{mdslate}{rgb}{0.45, 0.50, 0.68} % 114/127/173
\definecolor{ltslate}{rgb}{0.85, 0.88, 0.95} % 216/225/229


\newcommand{\includefigure}[2][]{{\centering\includegraphics[#1]{#2}\par}}
\newcommand{\highlight}[1]{{\bf\usebeamercolor[fg]{structure}#1}}
\newcommand{\important}[1]{{\color{red}#1}}
\newcommand{\issue}[2]{\item[Issue:] {\color{red}#1}\\{\item[Soln:] \color{blue}#2}\\[4pt]}

\setbeamercolor{alerted text}{fg=ltgreen}
\setbeamertemplate{description item}[align left]


\newcommand{\lhs}[1]{{\color{blue}#1}}
\newcommand{\rhs}[1]{{\color{red}#1}}
\newcommand{\annotateL}[2]{%
  {\color{blue}\underbrace{\color{blue}#1}_{\color{blue}\mathclap{#2}}}}
\newcommand{\annotateR}[2]{%
  {\color{red}\underbrace{\color{red}#1}_{\color{red}\mathclap{#2}}}}
\newcommand{\eqnannotate}[2]{%
  {\color{blue}%
  \underbrace{\color{black}#1}_{\color{blue}\mathclap{#2}}}}

\newcommand{\trialvec}[1][]{{\vec{\psi}_\mathit{trial}^{#1}}}
\newcommand{\trialscalar}[1][]{{\psi_\mathit{trial}^{#1}}}
\newcommand{\basisvec}[1][]{{\vec{\psi}_\mathit{basis}^{#1}}}
\newcommand{\basisscalar}[1][]{{\psi_\mathit{basis}^{#1}}}

\newcommand{\tensor}[1]{\bm{#1}}
\DeclareMathOperator{\Tr}{Tr}

\usefonttheme[onlymath]{serif}

% minted shortcuts
\newminted{cfg}{bgcolor=ltslate,autogobble,fontsize=\tiny}
\newminted{bash}{bgcolor=ltltgreen,autogobble,fontsize=\tiny}

% PyLith components
\newcommand{\pylith}[1]{{\ttfamily\color{magenta}#1}}



% ========================================================= DOCUMENT
\begin{document}

% ------------------------------------------------------------ SLIDE
\maketitle

\logo{\includegraphics[height=4.5ex]{../../logos/cig_short_pylith}}

% ========================================================== SECTION
\section{{\ttfamily examples/reverse-2d}}

% ========================================================== SUBSECTION
\subsection{Overview}

% ------------------------------------------------------------ SLIDE
\begin{frame}
  \frametitle{Vertical Cross-Section of Reverse Fault (2D): {\ttfamily examples/reverse-2d}}
  \summary{}

  \includefigure[height=6.1cm]{figs/geometry}

  \vfill
  Solve the static and quasistatic boundary elasticity equation for a vertical cross-section of a reverse fault with a splay.
  
\end{frame}


% ------------------------------------------------------------ SLIDE
\begin{frame}
  \frametitle{Steps in example}
  \summary{Steps 1-3 are covered in the gravitational body forces tutorial}

  \begin{description}
    \item[Step 1] Gravitational body forces and linear isotropic elasticity
    \item[Step 2] Gravitational body forces and linear isotropic elasticity with a reference stress state
    \item[Step 3] Gravitational body forces and linear isotropic incompressible elasticity
    \item[Step 4] Surface tractions and linear isotropic linear elasticity
    \item[Step 5] \highlight{Earthquake rupture on one fault and linear isotropic linear elasticity}
    \item[Step 6] \highlight{Earthquake rupture on two faults and linear isotropic linear elasticity}
    \item[Step 7] \highlight{Same as Step 6 but time-dependent with linear isotropic Maxwell viscoelastic rheology}
    \item[Step 8] Same as Step 7 but with linear isotropic power-law viscoelastic rheology
  \end{description}
  
\end{frame}


% ------------------------------------------------------------ SLIDE
\begin{frame}
  \frametitle{Concepts covered}
  \summary{}

  \begin{itemize}
  \item Generation of mesh using Gmsh and Cubit
  \item Variable mesh size with distance from the fault
  \item Static and quasi-static simulations for elasticity
  \item Prescribed slip earthquake rupture in 2D
  \item Prescribed slip on multiple faults
  \item Elastic and viscoelastic bulk rheologies
  \end{itemize}
  
\end{frame}

% ========================================================== SUBSECTION
\subsection{Finite-Element Mesh}

% ------------------------------------------------------------ SLIDE
\begin{frame}
  \frametitle{Gmsh: Geometry}
  \summary{}

  We construct the geometry by first creating points, then connecting the points into curves, and finally the curves into surfaces.
  
  \includefigure[height=6.5cm]{figs/geometry-gmsh}
  
\end{frame}


% ------------------------------------------------------------ SLIDE
\begin{frame}
  \frametitle{Gmsh: Creating the finite-element mesh}
  \summary{}

  \begin{onlyenv}<1>
    \begin{itemize}
    \item Each curve in Gmsh has a direction (orientation).
    \item The direction is from the starting point to the ending point.
    \item When connecting curves into surfaces, you must connect the curves in a consistent direction.
    \item We connect the curves in a counter-clockwise direction.
    \item To reverse the direction of a curve, use the negative tag.
    \end{itemize}
  \end{onlyenv}
  \begin{onlyenv}<2>
  \includefigure[height=7.0cm]{figs/gmsh-tri}
  \end{onlyenv}
  
\end{frame}


% ------------------------------------------------------------ SLIDE
\begin{frame}
  \frametitle{Cubit: Geometry}
  \summary{{\bfseries Note}: The provided parameter files use the Gmsh file}

  We construct the geometry by using geometric primitives and splitting surfaces using points and curves.
  
  \includefigure[height=6.5cm]{figs/geometry-cubit}
  
\end{frame}


% ------------------------------------------------------------ SLIDE
\begin{frame}
  \frametitle{Cubit: Creating the finite-element mesh}
  \summary{{\bfseries Note}: The provided parameter files use the Gmsh file}

  \includefigure[height=7.0cm]{figs/cubit-quad}
  
\end{frame}


% ========================================================== SECTION
\subsection{Files used for simulations}

% ------------------------------------------------------------ SLIDE
\begin{frame}
  \frametitle{Files used in simulations}
  \summary{Files are in directory {\tt examples/strikeslip-2d}}

  \begin{description}
  \item[README.md] Brief description of the various examples
  \item[*.cfg] PyLith parameter files
  \item[generate\_gmsh.py] Python script to generate mesh using Gmsh
  \item[*.msh] Finite-element mesh files generated by Gmsh
  \item[*.spatialdb] Spatial database files
  \item[viz] Directory containing ParaView Python scripts
  \item[output] Directory containing simulation output; created automatically when running the simulations
  \end{description}

\end{frame}


% ========================================================== SECTION
\subsection{step05-onefault}

% ------------------------------------------------------------ SLIDE
\begin{frame}
  \frametitle{Step 5: Overview}
  \summary{Uniform coseismic slip with Dirichlet (displacement) boundary conditions}

  \includefigure[height=6.5cm]{figs/step05-diagram}
      
\end{frame}


% ------------------------------------------------------------ SLIDE
\begin{frame}
  \frametitle{Step 5: Physics}
  \summary{}

  \begin{minipage}{0.3\textwidth}
    {\scriptsize
    \begin{gather*}
    % Solution
    \vec{s}^T = \left( \vec{u} \quad \vec{\lambda} \right)^T \\
    % Elasticity
    \tensor{\nabla} \cdot \tensor{\sigma}(\vec{u}) = \vec{0} \\
    % Dirichlet
    u_x = 0 \text{ on boundary\_xneg} \\
    u_x = 0 \text{ on boundary\_xpos} \\
    u_y = 0 \text{ on boundary\_yneg} \\
    % Prescribed slip
    d = d(y) \text{ on fault}
    \end{gather*}}
  \end{minipage}
  \hfill
  \begin{minipage}{0.67\textwidth}
    \includefigure[width=\textwidth]{figs/step05-diagram}
  \end{minipage}
      
\end{frame}


% ------------------------------------------------------------ SLIDE
\begin{frame}[t,fragile]
  \frametitle{Step 5: Physics to simulation parameters}
  \summary{}

  \begin{minipage}[t]{0.3\textwidth}
    {\scriptsize
    \begin{gather*}
    % Solution
    % Solution
    \vec{s}^T = \left( \vec{u} \quad \vec{\lambda} \right)^T \tikzmark{solution5}\\
    % Elasticity
    \tensor{\nabla} \cdot \tensor{\sigma}(\vec{u}) = \vec{0} \tikzmark{material5}\\
    % Dirichlet
    u_x = 0 \text{ on boundary\_xneg} \tikzmark{bc5}\\
    u_x = 0 \text{ on boundary\_xpos} \\
    u_y = 0 \text{ on boundary\_yneg} \\
    % Prescribed slip
    d = d(y) \text{ on fault} \tikzmark{fault5}
    \end{gather*}}
  \end{minipage}
  \hfill
  \begin{minipage}[t]{0.67\textwidth}
    % Solution
    \begin{onlyenv}<2>
      \tikzmark{solution5-cfg}
      \begin{cfgcode}
        [pylithapp.problem]
        solution = pylith.problems.SolnDispLagrange
        defaults.quadrature_order = 1
        
        [pylithapp.problem.solution.subfields]
        displacement.basis_order = 1
        lagrange_fault.basis_order = 1
      \end{cfgcode}
    \end{onlyenv}
    %
    % Governing equations
    \begin{onlyenv}<3>
      \tikzmark{material5-cfg}
      \begin{cfgcode}
        [pylithapp.problem]
        materials = [slab, plate, wedge]

        [pylithapp.problem.materials.slab]
        description = Slab material below main fault
        label_value = 1

        db_auxiliary_field = spatialdata.spatialdb.SimpleDB
        db_auxiliary_field.description = Elastic properties for slab
        db_auxiliary_field.iohandler.filename = mat_elastic.spatialdb

        auxiliary_subfields.density.basis_order = 0
        bulk_rheology.auxiliary_subfields.bulk_modulus.basis_order = 0
        bulk_rheology.auxiliary_subfields.shear_modulus.basis_order = 0

        derived_subfields.cauchy_strain.basis_order = 1
        derived_subfields.cauchy_stress.basis_order = 1
      \end{cfgcode}
    \end{onlyenv}
    %
    % Boundary conditions
    \begin{onlyenv}<4>
      \tikzmark{bc5-cfg}
      \begin{cfgcode}
        bc = [bc_xneg, bc_xpos, bc_yneg]
        bc.bc_xneg = pylith.bc.DirichletTimeDependent
        bc.bc_xpos = pylith.bc.DirichletTimeDependent
        bc.bc_yneg = pylith.bc.DirichletTimeDependent
        
        [pylithapp.problem.bc.bc_xpos]
        label = boundary_xpos
        label_value = 11
        constrained_dof = [0]
        db_auxiliary_field = pylith.bc.ZeroDB
        db_auxiliary_field.description = Dirichlet BC +x edge

        auxiliary_subfields.initial_amplitude.basis_order = 0 
      \end{cfgcode}
    \end{onlyenv}
    %
    % Fault
    \begin{onlyenv}<5>
      \tikzmark{fault5-cfg}
      \begin{cfgcode}
        [pylithapp.problem]
        interfaces = [fault]

        [pylithapp.problem.interfaces.fault]
        label = fault
        label_value = 20
        edge = fault_end
        edge_value = 21
        observers.observer.data_fields = [slip]

        [pylithapp.problem.interfaces.fault.eq_ruptures.rupture]
        db_auxiliary_field = spatialdata.spatialdb.UniformDB
        db_auxiliary_field.description = Fault rupture auxiliary field spatial database
        db_auxiliary_field.values = [initiation_time, final_slip_left_lateral, final_slip_opening]
        db_auxiliary_field.data = [0.0*s, -2.0*m, 0.0*m]
      \end{cfgcode}
    \end{onlyenv}
  \end{minipage}

  \begin{tikzpicture}[overlay,remember picture]
    \draw[physics-arrow,visible on=<2>] ($(pic cs:solution5-cfg)-(0,2em)$) to (pic cs:solution5);
    \draw[physics-arrow,visible on=<3>] ($(pic cs:material5-cfg)-(0,2em)$) to (pic cs:material5);
    \draw[physics-arrow,visible on=<4>] ($(pic cs:bc5-cfg)-(0,2em)$) to (pic cs:bc5);
    \draw[physics-arrow,visible on=<5>] ($(pic cs:fault5-cfg)-(0,2em)$) to (pic cs:fault5);
  \end{tikzpicture}
  
\end{frame}


% ------------------------------------------------------------ SLIDE
\begin{frame}
  \frametitle{Step 5: Input files}
  \summary{}

  \begin{description}
  \item[mesh\_tri.msh] Finite-element mesh generated using Gmsh
  \item[pylithapp.cfg] PyLith parameter file common to all steps
  \item[step05\_onefault.cfg] PyLith parameter file
  \item[mat\_elastic.spatialdb] Spatial database for isotropic linear elastic properties
  \end{description}
    
\end{frame}


% ------------------------------------------------------------ SLIDE
\begin{frame}[fragile]
  \frametitle{Step 5: Run the simulation}
  \summary{}

\begin{bashcode}
pylith step05_onefault.cfg

# Output
 >> /software/unix/py39-venv/pylith-debug/lib/python3.9/site-packages/pylith/meshio/MeshIOObj.py:44:read
 -- meshiopetsc(info)
 -- Reading finite-element mesh
 >> /src/cig/pylith/libsrc/pylith/meshio/MeshIO.cc:94:void pylith::meshio::MeshIO::read(pylith::topology::Mesh*)
 -- meshiopetsc(info)
 -- Component 'reader': Domain bounding box:
    (-100000, 100000)
    (-100000, 0)

# -- many lines omitted --

 >> /software/unix/py39-venv/pylith-debug/lib/python3.9/site-packages/pylith/problems/TimeDependent.py:139:run
 -- timedependent(info)
 -- Solving problem.
0 TS dt 0.01 time 0.
    0 SNES Function norm 2.129295960330e-02
    Linear solve converged due to CONVERGED_ATOL iterations 43
    1 SNES Function norm 1.577267844919e-12
  Nonlinear solve converged due to CONVERGED_FNORM_ABS iterations 1
1 TS dt 0.01 time 0.01
 >> /software/unix/py39-venv/pylith-debug/lib/python3.9/site-packages/pylith/problems/Problem.py:201:finalize
 -- timedependent(info)
 -- Finalizing problem.
\end{bashcode}
  
\end{frame}


% ------------------------------------------------------------ SLIDE
\begin{frame}
  \frametitle{Step 5: Visualize results}
  \summary{Run the {\tt viz/plot\_dispwarp.py} Python script from within ParaView}

  \includefigure[height=7.0cm]{figs/step05-solution}
    
\end{frame}


% ========================================================== SECTION
\subsection{step06-twofaults-elastic}

% ------------------------------------------------------------ SLIDE
\begin{frame}
  \frametitle{Step 6: Overview}
  \summary{Uniform rupture on two faults with Dirichlet (displacement) boundary conditions}

  \includefigure[height=6.5cm]{figs/step06-diagram}
      
\end{frame}


% ------------------------------------------------------------ SLIDE
\begin{frame}
  \frametitle{Step 6: Physics}
  \summary{}

  \begin{minipage}{0.3\textwidth}
    {\scriptsize
    \begin{gather*}
    % Solution
    \vec{s}^T = \left( \vec{u} \quad \vec{\lambda} \right)^T \\
    % Elasticity
    \tensor{\nabla} \cdot \tensor{\sigma}(\vec{u}) = \vec{0} \\
    % Dirichlet
    u_x = 0 \text{ on boundary\_xneg} \\
    u_x = 0 \text{ on boundary\_xpos} \\
    u_y = 0 \text{ on boundary\_yneg} \\
    % Prescribed slip
    d = d(y) \text{ on fault} \\
    d = d(y) \text{ on splay}
    \end{gather*}}
  \end{minipage}
  \hfill
  \begin{minipage}{0.67\textwidth}
    \includefigure[width=\textwidth]{figs/step06-diagram}
  \end{minipage}
      
\end{frame}


% ------------------------------------------------------------ SLIDE
\begin{frame}[t,fragile]
  \frametitle{Step 6: Physics to simulation parameters}
  \summary{}

  \begin{minipage}[t]{0.3\textwidth}
    {\scriptsize
    \begin{gather*}
    % Solution
    % Solution
    \vec{s}^T = \left( \vec{u} \quad \vec{\lambda} \right)^T \tikzmark{solution6}\\
    % Elasticity
    \tensor{\nabla} \cdot \tensor{\sigma}(\vec{u}) = \vec{0} \tikzmark{material6}\\
    % Dirichlet
    u_x = 0 \text{ on boundary\_xneg} \tikzmark{bc6}\\
    u_x = 0 \text{ on boundary\_xpos} \\
    u_y = 0 \text{ on boundary\_yneg} \\
    % Prescribed slip
    d = d(y) \text{ on fault} \tikzmark{fault6} \\
    d = d(y) \text{ on splay}
    \end{gather*}}
  \end{minipage}
  \hfill
  \begin{minipage}[t]{0.67\textwidth}
    % Solution
    \begin{onlyenv}<2>
      \tikzmark{solution6-cfg}
      \begin{cfgcode}
        # Same as Step 5
      \end{cfgcode}
    \end{onlyenv}
    %
    % Governing equations
    \begin{onlyenv}<3>
      \tikzmark{material6-cfg}
      \begin{cfgcode}
        # Same as Step 5
      \end{cfgcode}
    \end{onlyenv}
    %
    % Boundary conditions
    \begin{onlyenv}<4>
      \tikzmark{bc6-cfg}
      \begin{cfgcode}
        # Same as Step 5
      \end{cfgcode}
    \end{onlyenv}
    %
    % Fault
    \begin{onlyenv}<5>
      \tikzmark{fault6-cfg}
      \begin{cfgcode}
        [pylithapp.problem]
        interfaces = [fault, splay]

        [pylithapp.problem.interfaces.fault]
        label = fault
        label_value = 20
        edge = fault_end
        edge_value = 21
        observers.observer.data_fields = [slip]

        [pylithapp.problem.interfaces.fault.eq_ruptures.rupture]
        db_auxiliary_field = spatialdata.spatialdb.UniformDB
        db_auxiliary_field.description = Fault rupture auxiliary field spatial database
        db_auxiliary_field.values = [initiation_time, final_slip_left_lateral, final_slip_opening]
        db_auxiliary_field.data = [0.0*s, -2.0*m, 0.0*m]
      \end{cfgcode}
    \end{onlyenv}
    %
    % PETSc
    \begin{onlyenv}<6>
      \tikzmark{petsc6-cfg}
      \begin{cfgcode}
        # Wedge introduces null space; use ILU preconditioner
        [pylithapp.petsc]
        fieldsplit_displacement_pc_type = ilu
      \end{cfgcode}
    \end{onlyenv}
  \end{minipage}

  \begin{tikzpicture}[overlay,remember picture]
    \draw[physics-arrow,visible on=<2>] ($(pic cs:solution6-cfg)-(0,2em)$) to (pic cs:solution6);
    \draw[physics-arrow,visible on=<3>] ($(pic cs:material6-cfg)-(0,2em)$) to (pic cs:material6);
    \draw[physics-arrow,visible on=<4>] ($(pic cs:bc6-cfg)-(0,2em)$) to (pic cs:bc6);
    \draw[physics-arrow,visible on=<5>] ($(pic cs:fault6-cfg)-(0,2em)$) to (pic cs:fault6);
    \draw[physics-arrow,visible on=<6>] ($(pic cs:petsc6-cfg)-(0,2em)$) to (pic cs:fault6);
  \end{tikzpicture}
  
\end{frame}


% ------------------------------------------------------------ SLIDE
\begin{frame}
  \frametitle{Step 6: Input files}
  \summary{}

  \begin{description}
  \item[mesh\_tri.msh] Finite-element mesh generated using Gmsh
  \item[pylithapp.cfg] PyLith parameter file common to all steps
  \item[step06\_twofaults\_elastic.cfg] PyLith parameter file
  \item[mat\_elastic.spatialdb] Spatial database for isotropic linear elastic properties
  \end{description}
    
\end{frame}


% ------------------------------------------------------------ SLIDE
\begin{frame}[fragile]
  \frametitle{Step 6: Run the simulation}
  \summary{}

\begin{bashcode}
pylith step06_twofaults_elastic.cfg

# Output
 >> /software/unix/py39-venv/pylith-debug/lib/python3.9/site-packages/pylith/meshio/MeshIOObj.py:44:read
 -- meshiopetsc(info)
 -- Reading finite-element mesh
 >> /src/cig/pylith/libsrc/pylith/meshio/MeshIO.cc:94:void pylith::meshio::MeshIO::read(pylith::topology::Mesh*)
 -- meshiopetsc(info)
 -- Component 'reader': Domain bounding box:
    (-100000, 100000)
    (-100000, 0)

# -- many lines omitted --

 >> /software/unix/py39-venv/pylith-debug/lib/python3.9/site-packages/pylith/problems/TimeDependent.py:139:run
 -- timedependent(info)
 -- Solving problem.
0 TS dt 0.01 time 0.
    0 SNES Function norm 2.225574998436e-02
    Linear solve converged due to CONVERGED_ATOL iterations 415
    1 SNES Function norm 8.564274482242e-13
  Nonlinear solve converged due to CONVERGED_FNORM_ABS iterations 1
1 TS dt 0.01 time 0.01
 >> /software/unix/py39-venv/pylith-debug/lib/python3.9/site-packages/pylith/problems/Problem.py:201:finalize
 -- timedependent(info)
 -- Finalizing problem.
\end{bashcode}
  
\end{frame}


% ------------------------------------------------------------ SLIDE
\begin{frame}
  \frametitle{Step 6: Visualize results}
  \summary{Run the {\tt viz/plot\_dispwarp.py} Python script from within ParaView}

  \includefigure[height=7.0cm]{figs/step06-solution}
    
\end{frame}


% ========================================================== SECTION
\subsection{step07-twofaults-maxwell}

% ------------------------------------------------------------ SLIDE
\begin{frame}
  \frametitle{Step 7: Overview}
  \summary{Uniform rupture on two faults with Dirichlet (displacement) boundary conditions}

  Same as Step 6 with linear Maxwell viscoelastic bulk rheology for the slab.
      
\end{frame}


% ------------------------------------------------------------ SLIDE
\begin{frame}[t,fragile]
  \frametitle{Step 7: Physics to simulation parameters}
  \summary{}

  \begin{minipage}[t]{0.3\textwidth}
    {\scriptsize
    \begin{gather*}
    % Solution
    % Solution
    \vec{s}^T = \left( \vec{u} \quad \vec{\lambda} \right)^T \tikzmark{solution7}\\
    % Elasticity
    \tensor{\nabla} \cdot \tensor{\sigma}(\vec{u}) = \vec{0} \tikzmark{material7}\\
    \rho, v_p, v_s \text{ in crust and wedge} \\
    \rho, v_p, v_s, \nu  \text{ in slab} \\
    % Dirichlet
    u_x = 0 \text{ on boundary\_xneg} \tikzmark{bc7}\\
    u_x = 0 \text{ on boundary\_xpos} \\
    u_y = 0 \text{ on boundary\_yneg} \\
    % Prescribed slip
    d = d(y) \text{ on fault} \tikzmark{fault7} \\
    d = d(y) \text{ on splay}
    \end{gather*}}
  \end{minipage}
  \hfill
  \begin{minipage}[t]{0.67\textwidth}
    % Solution
    \begin{onlyenv}<2>
      \tikzmark{solution7-cfg}
      \begin{cfgcode}
        [pylithapp.problem]
        initial_dt = 4.0*year
        start_time = -4.0*year
        end_time = 100.0*year

        normalizer.relaxation_time = 20.0*year
      \end{cfgcode}
    \end{onlyenv}
    %
    % Governing equations
    \begin{onlyenv}<3>
      \tikzmark{material7-cfg}
      \begin{cfgcode}
        [pylithapp.problem.materials.slab]
        description = Slab material below main fault
        label_value = 1
        bulk_rheology = pylith.materials.IsotropicLinearMaxwell

        db_auxiliary_field = spatialdata.spatialdb.SimpleDB
        db_auxiliary_field.description = Elastic properties for slab
        db_auxiliary_field.iohandler.filename = mat_maxwell.spatialdb

        auxiliary_subfields.density.basis_order = 0
        bulk_rheology.auxiliary_subfields.bulk_modulus.basis_order = 0
        bulk_rheology.auxiliary_subfields.shear_modulus.basis_order = 0
        bulk_rheology.auxiliary_subfields.maxwell_time.basis_order = 0
        
        derived_subfields.cauchy_strain.basis_order = 1
        derived_subfields.cauchy_stress.basis_order = 1
      \end{cfgcode}
    \end{onlyenv}
    %
    % Boundary conditions
    \begin{onlyenv}<4>
      \tikzmark{bc7-cfg}
      \begin{cfgcode}
        # Same as Step 6
      \end{cfgcode}
    \end{onlyenv}
    %
    % Fault
    \begin{onlyenv}<5>
      \tikzmark{fault7-cfg}
      \begin{cfgcode}
        # Same as Step 6
      \end{cfgcode}
    \end{onlyenv}
    %
    % PETSc
    \begin{onlyenv}<6>
      \tikzmark{petsc7-cfg}
      \begin{cfgcode}
        # Wedge introduces null space; use ILU preconditioner
        # Same as Step 6
        [pylithapp.petsc]
        fieldsplit_displacement_pc_type = ilu
      \end{cfgcode}
    \end{onlyenv}
  \end{minipage}

  \begin{tikzpicture}[overlay,remember picture]
    \draw[physics-arrow,visible on=<2>] ($(pic cs:solution7-cfg)-(0,2em)$) to (pic cs:solution7);
    \draw[physics-arrow,visible on=<3>] ($(pic cs:material7-cfg)-(0,2em)$) to (pic cs:material7);
    \draw[physics-arrow,visible on=<4>] ($(pic cs:bc7-cfg)-(0,2em)$) to (pic cs:bc7);
    \draw[physics-arrow,visible on=<5>] ($(pic cs:fault7-cfg)-(0,2em)$) to (pic cs:fault7);
    \draw[physics-arrow,visible on=<6>] ($(pic cs:petsc7-cfg)-(0,2em)$) to (pic cs:fault7);
  \end{tikzpicture}
  
\end{frame}


% ------------------------------------------------------------ SLIDE
\begin{frame}
  \frametitle{Step 7: Input files}
  \summary{}

  \begin{description}
  \item[mesh\_tri.msh] Finite-element mesh generated using Gmsh
  \item[pylithapp.cfg] PyLith parameter file common to all steps
  \item[step07\_twofaults\_maxwell.cfg] PyLith parameter file
  \item[mat\_elastic.spatialdb] Spatial database for isotropic linear elastic properties
  \item[mat\_maxwell.spatialdb] Spatial database for isotropic linear Maxwell properties
  \end{description}
    
\end{frame}


% ------------------------------------------------------------ SLIDE
\begin{frame}[fragile]
  \frametitle{Step 7: Run the simulation}
  \summary{}

\begin{bashcode}
pylith step07_twofaults_maxwell.cfg

# Output
 >> /software/unix/py39-venv/pylith-debug/lib/python3.9/site-packages/pylith/meshio/MeshIOObj.py:44:read
 -- meshiopetsc(info)
 -- Reading finite-element mesh
 >> /src/cig/pylith/libsrc/pylith/meshio/MeshIO.cc:94:void pylith::meshio::MeshIO::read(pylith::topology::Mesh*)
 -- meshiopetsc(info)
 -- Component 'reader': Domain bounding box:
    (-100000, 100000)
    (-100000, 0)

# -- many lines omitted --

25 TS dt 0.2 time 4.8
    0 SNES Function norm 2.589196279152e-05
    Linear solve converged due to CONVERGED_ATOL iterations 339
    1 SNES Function norm 6.643244443329e-13
  Nonlinear solve converged due to CONVERGED_FNORM_ABS iterations 1
26 TS dt 0.2 time 5.
 >> /software/unix/py39-venv/pylith-debug/lib/python3.9/site-packages/pylith/problems/Problem.py:201:finalize
 -- timedependent(info)
 -- Finalizing problem.
\end{bashcode}
  
\end{frame}


% ------------------------------------------------------------ SLIDE
\begin{frame}
  \frametitle{Step 7: Visualize results}
  \summary{Run the {\tt viz/plot\_dispwarp.py} Python script from within ParaView}

  \includefigure[height=7.0cm]{figs/step07-solution}
    
\end{frame}



% ======================================================================
\end{document}


% End of file
